\documentclass[conference]{IEEEtran}

\usepackage[T1]{fontenc}
\usepackage[utf8]{inputenc}
\usepackage{cite}
\usepackage{amsmath}
\usepackage{graphicx}
% Extra packages for figures and captions
\usepackage{caption}
\usepackage{subcaption}
% Look for images in these folders (place your images here)
\graphicspath{{images/}{figures/}}
\usepackage{booktabs}
\usepackage{array}
\usepackage{url}
\usepackage{xcolor}
% Hyperlinks (load near end of preamble)
\usepackage[hidelinks]{hyperref}

\title{DevRoom: A Real-Time Collaborative Workspace for Software Architecture Design and File-Centric Team Coordination}

\author{
\IEEEauthorblockN{Author Name}
\IEEEauthorblockA{
Department or Institution Name \\
City, Country \\
email@example.com
}
}

\begin{document}

\maketitle

\begin{abstract}
Modern software development workflows distribute design, documentation, and code across multiple specialized tools. This fragmentation increases context switching and reduces traceability between high-level architecture decisions and implementation artifacts. DevRoom is a prototype full-stack workspace that integrates real-time architecture diagramming, file versioning, node-to-file linking, role-aware room membership, and a persistent activity history. Implemented with a React + Vite frontend, a Node.js + Express backend, MongoDB persistence, and Socket.IO for synchronization, DevRoom demonstrates a compact, extensible architecture for bridging visual design and file-centric collaboration. This paper describes the design and implementation of DevRoom, highlights engineering trade-offs, and outlines directions for evaluation and production hardening.
\end{abstract}

\begin{IEEEkeywords}
collaborative systems, software architecture, real-time collaboration, React, Node.js, MongoDB, Socket.IO, file versioning
\end{IEEEkeywords}

\section{Introduction}
Teams building modern distributed systems must reason about architecture, interfaces, and evolving code simultaneously. Designers often sketch high-level architectures in whiteboard-style tools while the corresponding implementation and documentation live in separate repositories and editors. This separation increases cognitive overhead and complicates traceability from design elements to concrete artifacts.

DevRoom aims to reduce these gaps by co-locating diagrammatic design and file-centric artifacts inside a room-based collaboration environment. The prototype supports synchronous editing of a shared visual canvas, links between canvas nodes and file versions, membership and role controls for rooms, and a persistent activity log that captures key events. The system is intentionally lightweight and modular to make integration into existing workflows straightforward.

The main contributions of this work are:
\begin{itemize}
\item A description of an integrated workspace that unifies real-time diagramming with file versioning and node-to-file links.
\item A full-stack prototype demonstrating practical implementation choices and trade-offs for collaborative design tools.
\item A qualitative discussion of limitations and a roadmap for empirical evaluation and production readiness.
\end{itemize}

\section{Related Work}
Collaborative editing and whiteboarding tools have been widely studied and deployed. Real-time document editors (e.g., Google Docs) and drawing tools (e.g., Miro, Figma) provide low-latency shared editing and presence. Tooling focused on architecture and modelling (e.g., PlantUML, diagramming platforms) emphasizes expressive visual notation but often lacks tight integration with code and versioned artifacts. Research on traceability and software archaeology explores connections between design rationale and implementation history. DevRoom builds on these ideas by focusing specifically on linking visual architecture nodes to versioned files and preserving a room-level activity history.

\section{Problem Statement}
Current developer workflows commonly show the following issues:
\begin{itemize}
\item Design artifacts and implementation assets are kept in separate tools, reducing traceability.
\item Realtime collaboration often focuses on drawing or text editing in isolation rather than connecting artifacts.
\item Access control and membership are handled outside the design environment, complicating coordination.
\item Changes across artifacts lack a unified activity history for review and auditing.
\end{itemize}

DevRoom targets these gaps by offering a single workspace where diagrams, files, membership, and history coexist.

\section{Objectives}
The prototype is guided by the following objectives:
\begin{itemize}
\item Provide a shared visual canvas for architecture planning and quick exploration.
\item Support real-time collaboration (presence, cursors, concurrent edits).
\item Allow users to upload, edit, and version text-oriented project files.
\item Enable direct links from canvas nodes to specific files and file versions.
\item Manage rooms and collaborators with role-based permissions.
\item Record an activity history for auditing and asynchronous review.
\end{itemize}

\section{System Architecture}
DevRoom follows a classical client-server architecture augmented with a realtime synchronization layer. The frontend provides an interactive UI and the canvas editor; the backend exposes RESTful APIs for persistence and business logic. Socket.IO provides low-latency synchronization for presence and immediate canvas events.

\subsection{Frontend Layer}
The frontend is implemented with React and Vite. Key components include:
\begin{itemize}
\item Authentication and onboarding screens.
\item Room lifecycle screens (creation, listing, joining).
\item Main workspace with tabs for Canvas, Files, Members, and History.
\item File editor modal that supports editing and version commits.
\end{itemize}

React Flow is used to implement the architecture canvas because it provides a flexible node/edge model and integrates well with React state and events.
\begin{figure}[!t]
	\centering
	\includegraphics[width=\columnwidth]{architecture.png}
	\caption{DevRoom system architecture (frontend, realtime layer, backend, persistence).}
	\label{fig:architecture}
\end{figure}

\subsection{Backend Layer}
The backend uses Node.js and Express with Mongoose for MongoDB access. Primary responsibilities include user management, room and membership handling, file storage and versioning metadata, canvas persistence, and history logging. JWT tokens protect API endpoints and role checks are enforced by middleware.

\subsection{Realtime Layer}
Socket.IO handles presence, cursors, and lightweight canvas events (node creation, selection, transient movements). For durability, long-lived changes (full canvas saves, file commits) are persisted via REST endpoints; realtime messages are used to propagate ephemeral state quickly between clients.

\section{Data Model}
The persistence model focuses on a compact set of collections that capture users, rooms, canvas elements, files, and history. Table~\ref{tab:data-model} summarizes the principal entities.

\begin{table}[htbp]
\caption{Core Data Entities in DevRoom}
\label{tab:data-model}
\centering
\begin{tabular}{>{\raggedright\arraybackslash}p{2.2cm} >{\raggedright\arraybackslash}p{5.0cm}}
\toprule
Entity & Purpose \\
\midrule
User & Registered users and authentication profile data \\
Room & Represents a collaborative workspace (room-level settings, metadata) \\
RoomMember & Maps users to rooms and stores roles (owner, editor, viewer) \\
CanvasNode & Diagram nodes with type, title, properties, and position \\
CanvasEdge & Directed links between nodes and associated metadata \\
CanvasNote & Sticky-note annotations attached to the canvas \\
File & Stores the current file snapshot and basic metadata \\
FileVersion & Preserves previous file revisions with commit notes \\
FileNodeLink & Associates files or specific versions with canvas nodes \\
HistoryLog & Time-ordered entries describing room events and actor information \\
\bottomrule
\end{tabular}
\end{table}

\begin{figure}[!t]
	\centering
	\includegraphics[width=\columnwidth]{data-model.png}
	\caption{Simplified data model for DevRoom (collections and relations).}
	\label{fig:data-model}
\end{figure}

\section{Implementation Details}

\subsection{Authentication and Authorization}
DevRoom uses JWT-based authentication and bcrypt for password hashing. Tokens are issued on successful login and attached to protected requests. Role checks are enforced at the route level using middleware that inspects the \texttt{RoomMember} mapping. The prototype stores tokens in browser local storage for session continuity; production deployments should consider secure cookie storage and CSRF protections.

\subsection{Collaborative Canvas}
The canvas supports several predefined node types (e.g., service, database, gateway) and free-form sticky notes. Client-side edits are emitted over Socket.IO for immediate feedback. Durable changes can be persisted incrementally or via a full-canvas save operation. The latter replaces the canonical canvas snapshot in the database and reconciles temporary frontend identifiers with server-assigned MongoDB IDs to maintain stable references.

\subsection{File Management}
DevRoom currently targets text-oriented files (design notes, small snippets, configuration). Each file has a canonical current snapshot recorded in the \texttt{File} collection; historical revisions are stored in \texttt{FileVersion} with optional commit messages. Files may be linked to canvas nodes using \texttt{FileNodeLink}, enabling quick navigation from design elements to their corresponding artifacts.

\subsection{Activity History}
The backend records history entries for important actions (node/edge creation, file commits, linking operations, membership changes). The History view exposes these entries in reverse chronological order and includes metadata such as actor, timestamp, and optional human-readable notes to support asynchronous review.

\section{Operational Workflow}
Typical usage follows these steps:
\begin{enumerate}
\item Register and authenticate with the system.
\item Create or join a room and obtain the appropriate role.
\item Collaborators open the shared canvas and add nodes, edges, and notes.
\item Upload or edit files and optionally link specific versions to nodes.
\item Use presence and cursors to coordinate live editing; persist important changes via commits or full-canvas saves.
\item Inspect the activity history for a chronological account of room changes.
\end{enumerate}

\section{Prototype Demonstration and Discussion}
The DevRoom prototype demonstrates several practical engineering trade-offs for integrated collaboration:
\begin{itemize}
\item Performance: using Socket.IO for ephemeral events keeps latency low for interactive feedback, while writes that must be durable are done through REST to simplify consistency semantics.
\item Traceability: file-to-node links provide a simple, navigable connection between design elements and artifacts, improving discoverability compared with separate tools.
\item Extensibility: the separation of realtime events and persistent APIs makes it feasible to add features such as background indexing, richer file types, or external storage backends.
\end{itemize}

\begin{figure}[!t]
	\centering
	\includegraphics[width=\columnwidth]{frontend-ui.png}
	\caption{Frontend workspace screenshot: canvas, files, and panels.}
	\label{fig:frontend-ui}
\end{figure}

The implementation is a prototype and not yet a full empirical study. We view the system as a platform for follow-up user studies and for iterating on conflict resolution and scalability strategies.

\section{Current Limitations}
The prototype exposes several limitations that should be addressed before wider deployment:
\begin{itemize}
\item Some file-related endpoints are not exhaustively guarded by authorization checks in the current codebase.
\item Configuration values (API and socket endpoints) are currently hard-coded for local development.
\item Realtime presence is stored in-memory on the server and is lost on backend restarts; a resilient presence store is needed for production.
\item Conflict handling for concurrent edits is rudimentary and requires stronger merging or CRDT-based approaches for richer file types.
\item The system currently targets text files; supporting binary or large media attachments requires design changes (external storage, streaming uploads).
\end{itemize}

\section{Future Work}
Priority next steps include:
\begin{itemize}
\item Externalize configuration into environment variables and secure deployment templates (Docker, CI pipelines).
\item Add comprehensive authorization checks and secure token handling.
\item Replace in-memory presence with a distributed store (Redis) to survive restarts.
\item Explore CRDTs or operational transform strategies to improve concurrent edit handling for both canvas state and file contents.
\item Conduct formal user studies to measure productivity, traceability improvements, and usability feedback.
\end{itemize}

\section{Conclusion}
DevRoom presents an integrated approach to combining visual software architecture design with file-centric collaboration and history tracking. The prototype demonstrates that a compact stack (React, Node.js, MongoDB, Socket.IO) can be composed to provide synchronous canvas editing, file versioning, and navigable links between design and implementation artifacts. Future work will focus on hardening, scalability, and empirical evaluation.

\section*{Acknowledgment}
This paper reflects the current prototype implementation of the DevRoom project. Please replace author and affiliation placeholders with actual details before formal submission.

\begin{thebibliography}{00}
\bibitem{reactflow}
React Flow, ``React Flow Documentation,'' Available: \url{https://reactflow.dev/}

\bibitem{react}
React, ``React Documentation,'' Available: \url{https://react.dev/}

\bibitem{nodejs}
Node.js, ``Node.js Documentation,'' Available: \url{https://nodejs.org/}

\bibitem{express}
Express, ``Express Documentation,'' Available: \url{https://expressjs.com/}

\bibitem{mongoose}
Mongoose, ``Mongoose Documentation,'' Available: \url{https://mongoosejs.com/}

\bibitem{socketio}
Socket.IO, ``Socket.IO Documentation,'' Available: \url{https://socket.io/docs/}

\bibitem{mongodb}
MongoDB, ``MongoDB Documentation,'' Available: \url{https://www.mongodb.com/docs/}
\end{thebibliography}
%
% -----------------------------------------------------------------------------
% IMAGE INSTRUCTIONS (place these files in an "images/" folder next to this .tex file)
% Required images (suggested names):
% - architecture.png      : System architecture diagram (frontend, socket layer, backend, DB)
% - data-model.png        : ER / collection diagram for the database model
% - frontend-ui.png       : Screenshot of the frontend workspace (canvas + panels)
% - canvas-screenshot.png : Example canvas showing nodes and edges (optional)
% Provide PNG or PDF images; 300 DPI recommended. Filenames must match exactly.
% -----------------------------------------------------------------------------

\end{document}
